\documentclass[aspectratio=169]{beamer}

\usepackage{graphicx}
\usepackage{booktabs}
\usepackage{amsmath}
\usepackage{amssymb}
\usepackage{pgfpages}
\usepackage{xspace}

\graphicspath{{figs/}}
\setbeameroption{show notes}
\addtobeamertemplate{note page}{}{\thispdfpagelabel{notes:\insertframenumber}}
\setbeamerfont{note page}{size=\footnotesize}
\makeatletter
\defbeamertemplate*{note page}{threelevel}
{%
  {%
    \footnotesize
    \usebeamerfont{note title}\usebeamercolor[fg]{note title}%
    \ifbeamercolorempty[bg]{note title}{}{%
      \insertvrule{.25\paperheight}{note title.bg}%
      \vskip-.25\paperheight%
      \nointerlineskip%
    }%
    \vbox{
      \hfill\insertslideintonotes{0.25}\hskip-\Gm@rmargin\hskip0pt%
      \vskip-0.25\paperheight%
      \nointerlineskip
      \begin{pgfpicture}{0cm}{0cm}{0cm}{0cm}
        \begin{pgflowlevelscope}{\pgftransformrotate{90}}
          {\pgftransformshift{\pgfpoint{-2cm}{0.2cm}}%
          \pgftext[base,left]{\usebeamerfont{note date}\usebeamercolor[fg]{note date}\the\year-\ifnum\month<10\relax0\fi\the\month-\ifnum\day<10\relax0\fi\the\day}}
        \end{pgflowlevelscope}
      \end{pgfpicture}}
    \nointerlineskip
    \vbox to .25\paperheight{\vskip0.5em
      \hbox{\insertshorttitle[width=0.75\textwidth]}%
      \setbox\beamer@tempbox=\hbox{\insertsection}%
      \hbox{\ifdim\wd\beamer@tempbox>1pt{\hskip4pt\raise3pt\hbox{\vrule
            width0.4pt height7pt\vrule width 9pt
            height0.4pt}}\hskip1pt\hbox{\begin{minipage}[t]{0.71\textwidth}\def\breakhere{}\insertsection\end{minipage}}\fi%
      }%
      \setbox\beamer@tempbox=\hbox{\insertsubsection}%
      \hbox{\ifdim\wd\beamer@tempbox>1pt{\hskip17.4pt\raise3pt\hbox{\vrule
            width0.4pt height7pt\vrule width 9pt
            height0.4pt}}\hskip1pt\hbox{\begin{minipage}[t]{0.71\textwidth}\def\breakhere{}\insertsubsection\end{minipage}}\fi%
      }%
      \vfil}%
  }%
  \ifbeamercolorempty[bg]{note page}{}{%
    \nointerlineskip%
    \insertvrule{.75\paperheight}{note page.bg}%
    \vskip-.75\paperheight%
  }%
  \vskip.25em
  \nointerlineskip
  \insertnote
}
\setbeamertemplate{note page}[threelevel]
\makeatother

% Show "current / total" page number in bottom-right; hide navigation icons
\setbeamertemplate{navigation symbols}{}
\setbeamertemplate{footline}{%
  \hbox{%
    \begin{beamercolorbox}[wd=\paperwidth,ht=2.25ex,dp=1ex,right]{}%
      \usebeamerfont{footline}\color{gray}\insertframenumber\,/\,\inserttotalframenumber\hspace*{2.5ex}%
    \end{beamercolorbox}%
  }%
}

\newcommand{\design}{\textit{FastTT}\xspace}
\newcommand{\sxor}{Shift-XOR\xspace}
\newcommand{\twotone}{Two-tone\xspace}

\title{\design: Accelerating \sxor Erasure Coding for Data Storage}
\author{Duo Sun \and Ximing Fu \and Qiang Wang}
\institute{Harbin Institute of Technology (Shenzhen) \\ Pengcheng Laboratory}
\date{IPDPS 2026}

\begin{document}

\part{\design: Accelerating \sxor Erasure Coding for Data Storage}

% =====================================================================
% Title
% =====================================================================
\begin{frame}[plain]
  \titlepage
\end{frame}
\note{
Good afternoon. I am Duo Sun, presenting FastTT: Accelerating Shift-XOR Erasure Coding for Data Storage, joint work with my supervisors Prof. Ximing Fu and Prof. Qiang Wang at HIT Shenzhen. Two-tone is a Shift-XOR erasure code that replaces finite-field multiplications with bit-shifts and XORs. This talk is about turning that theoretical efficiency into a practical, high-throughput implementation inside Ceph.
}

% =====================================================================
% Motivation
% =====================================================================
\section{Motivation}
\subsection{Why This Problem Matters}

\begin{frame}[t]{Motivation}
\small
\begin{itemize}
  \setlength{\itemsep}{0.55em}
  \item Erasure coding $=$ reliability + low storage cost.
  \item \textbf{Reed-Solomon (RS) Codes:} mature; many system-level optimizations already exist.
  \item \textbf{\twotone \sxor:} excellent theoretical complexity --- only bit-shifts \& XORs, no finite-field multiplications.
  \begin{itemize}
    \setlength{\itemsep}{0.25em}
    \item But a \emph{direct} implementation is \textbf{memory-bound}.
    \item RS optimizations do \textbf{not} transfer --- different compute \& memory-access patterns.
  \end{itemize}
\end{itemize}

\vspace{0.9em}
\begin{block}{Our Goal}
First high-performance, memory-access-aware implementation of \twotone \sxor, validated in Ceph.
\end{block}
\end{frame}
\note{
Erasure coding delivers reliability at much lower storage cost than replication. RS Codes are the dominant choice and already heavily optimized. Our focus is Two-tone Shift-XOR: theoretically excellent, only shifts and XORs, no finite-field multiplications. The gap is on the implementation side: a direct Two-tone implementation is memory-bound, and RS-specific tricks do not transfer because the computation and memory-access patterns are fundamentally different. FastTT is the first high-performance, memory-access-aware implementation, validated in Ceph.
}

% =====================================================================
% Background: Two-tone at a glance
% =====================================================================
% =====================================================================
% Limitations of existing solutions
% =====================================================================
\section{Limitations of Existing Solutions}
\subsection{Encoding \& Decoding Bottlenecks}

\begin{frame}[t]{Limitations of Existing Solutions}
\scriptsize
\begin{columns}[T,totalwidth=\textwidth]
  % --- Left: Encoding & Decoding (merged) ---
  \column{0.42\textwidth}
  \textbf{Encoding \& Decoding}
  \begin{itemize}
    \setlength{\itemsep}{0.25em}
    \item Algorithm \& correctness established in prior work --- briefly recalled here.
    \item Two traversal patterns of the data-parity grid: \emph{horizontal} (default direct implementation) and \emph{vertical}.
    \item \textbf{Decoding:} generic 3-stage flow --- Substitution $\to$ \twotone Elimination $\to$ Repair (see fig.~(c)).
  \end{itemize}

  % --- Right: traversal patterns + decoding example side-by-side ---
  \column{0.58\textwidth}
  \begin{minipage}[t]{0.32\linewidth}
    \vspace*{0pt}\centering
    \includegraphics[width=\linewidth,height=0.18\textheight,keepaspectratio]{Two-tone_Horizontal_Access_Pattern.pdf}\\[-0.2em]
    {\tiny (a) Horizontal traversal} \\[-0.3em]
    {\tiny \emph{(default encoding)}}

    \vspace{0.2em}
    \includegraphics[width=\linewidth,height=0.18\textheight,keepaspectratio]{Two-tone_Vertical_Access_Pattern.pdf}\\[-0.2em]
    {\tiny (b) Vertical traversal}
  \end{minipage}\hfill
  \begin{minipage}[t]{0.66\linewidth}
    \vspace*{0pt}\centering
    \includegraphics[width=\linewidth,height=0.43\textheight,keepaspectratio]{Example_of_Two-tone_Decoding.pdf}\\[-0.2em]
    {\tiny (c) \twotone decoding example: $d_1,d_2,c_2$ erased} \\[-0.3em]
    {\tiny Substitution $\to$ Elimination $\to$ Repair}
  \end{minipage}
\end{columns}

\vspace{0.05em}
\centering
\begin{minipage}{0.96\linewidth}
\begin{block}{Bottleneck $=$ memory access, not arithmetic}
\vspace{-0.4em}
\scriptsize
\textbf{Both encoding \& decoding:} poor locality, repeated access of the same chunks. \\
\textbf{Decoding extra:} recovering lost \emph{data} and re-encoding lost \emph{parity} revisit overlapping memory.
\vspace{-0.3em}
\end{block}
\end{minipage}
\end{frame}
\note{
The Two-tone encoding and decoding algorithms, together with their correctness, were established in prior work, so I will skip the math and focus on the implementation gap. For encoding, horizontal traversal of the data-parity grid is the default direct implementation, with vertical as the alternative. Decoding follows a generic three-stage flow --- substitute survived data, run Two-tone elimination, then repair the lost parity --- illustrated on the right with $d_1$, $d_2$ and $c_2$ erased. The bottleneck is the same on both sides: poor locality and repeated access over the same chunks. Decoding pays an extra cost: solving for lost data and re-encoding lost parity are separate passes that revisit overlapping memory. The root cause is memory access behavior, not coding arithmetic.
}

% =====================================================================
% Approach and key ideas
% =====================================================================
\section{Approach and Key Ideas}
\subsection{FastTT Overview}

\begin{frame}[t]{\design Overview}
\footnotesize
\begin{columns}[T,totalwidth=\textwidth]
  \column{0.42\textwidth}
  \begin{block}{Design principle}
  \small Optimize \twotone by \textbf{memory-access behavior}, not by algorithmic rewriting.
  \end{block}

  \vspace{0.3em}
  \textbf{Four coordinated techniques}
  \begin{enumerate}
    \setlength{\itemsep}{0.25em}
    \item \textbf{CFWP} \,-- Cache-Friendly Window Partitioning
    \item \textbf{NMA} \,-- Neighbor-Merging Access
    \item \textbf{EPM} \,-- Erasure-Parity Mapping
    \item \textbf{TPS} \,-- Traversal Pattern Selection
  \end{enumerate}

  \vspace{0.3em}
  \textit{Encoding:} CFWP $+$ NMA. \\
  \textit{Decoding:} CFWP $+$ NMA $+$ EPM $+$ TPS.

  \column{0.58\textwidth}
  \centering
  \includegraphics[width=\linewidth,height=0.7\textheight,keepaspectratio]{Two-tone_Optimization_Overview.pdf}

  \vspace{0.2em}
  {\scriptsize Overall \design workflow.}
\end{columns}
\end{frame}
\note{
FastTT follows one design principle. Instead of rewriting the Two-tone code construction, we optimize its memory-access behavior. Four coordinated techniques deliver this. Encoding uses Cache-Friendly Window Partitioning together with Neighbor-Merging Access to keep the working set in cache and cut redundant parity writes. Decoding reuses the same two and adds Erasure-Parity Mapping, which fuses data recovery with parity re-encoding into a single pass, and Traversal Pattern Selection, which picks the best access order for the erasure pattern at hand. The next four slides walk through these techniques one by one.
}

\subsection{Cache-Friendly Window Partitioning}

\begin{frame}[t]{Key Idea 1 -- Cache-Friendly Window Partitioning}
\footnotesize
\begin{columns}[T,totalwidth=\textwidth]
  \column{0.42\textwidth}
  \begin{itemize}
    \setlength{\itemsep}{0.45em}
    \item Tile each chunk into \textbf{cache-sized windows}.
    \item Finish all \sxor updates in one window before moving on.
    \item \textbf{On-demand init} -- no full zero pass up front.
  \end{itemize}

  \vspace{0.4em}
  \begin{block}{Effect}
  \small Localized access $+$ one fewer full traversal. \\
  Benefits both encoding and decoding.
  \end{block}

  \column{0.58\textwidth}
  \centering
  \includegraphics[width=\linewidth,height=0.65\textheight,keepaspectratio]{Cache-Friendly_Window_Demonstration.pdf}

  \vspace{0.1em}
  {\scriptsize Each chunk is tiled into windows; all parity updates inside a window finish before advancing. Colored segments show local access.}
\end{columns}
\end{frame}
\note{
Throughout the key ideas I will use the (4,3) code as a running example, matching the figures. The idea of Cache-Friendly Window Partitioning is to tile each chunk into windows that fit in cache, and finish all Shift-XOR updates inside a window before moving to the next. The working set stays small and fits in cache. On-demand initialization also removes one full zero-pass over the chunk. This benefits both encoding and decoding.
}

\subsection{Neighbor-Merging Access}

\begin{frame}[t]{Key Idea 2 -- Neighbor-Merging Access}
\footnotesize
\begin{columns}[T,totalwidth=\textwidth]
  \column{0.36\textwidth}
  \begin{itemize}
    \setlength{\itemsep}{0.5em}
    \item XOR \textbf{two adjacent} data symbols in SIMD registers \emph{first}.
    \item Then write the merged result to parity \emph{once}.
    \item Keeps horizontal-style read locality; halves parity writes.
  \end{itemize}

  \vspace{0.5em}
  \begin{block}{Effect}
  \small Fewer redundant accesses than \emph{both} horizontal and vertical traversals in practice.
  \end{block}

  \column{0.64\textwidth}
  \centering
  \begin{minipage}[t]{0.48\linewidth}
    \vspace*{0pt}\centering
    \includegraphics[width=\linewidth,height=0.26\textheight,keepaspectratio]{Two-tone_Horizontal_Access_Pattern.pdf}\\
    {\tiny (a) Horizontal}
  \end{minipage}\hfill
  \begin{minipage}[t]{0.48\linewidth}
    \vspace*{0pt}\centering
    \includegraphics[width=\linewidth,height=0.26\textheight,keepaspectratio]{Two-tone_Vertical_Access_Pattern.pdf}\\
    {\tiny (b) Vertical}
  \end{minipage}

  \vspace{0.5em}
  \begin{minipage}[t]{0.98\linewidth}
    \vspace*{0pt}\centering
    \includegraphics[width=\linewidth,height=0.42\textheight,keepaspectratio]{Neighbor-Merging_Access_Pattern.pdf}\\
    {\scriptsize (c) \textbf{Neighbor-Merging.} Adjacent data symbols are XORed in register, then written once. Parity writes: $4 \to 2$ per position.}
  \end{minipage}
\end{columns}
\end{frame}
\note{
In the example, two adjacent data chunks share the same offset to a particular parity, so each contributes to that parity slot without any shift. Instead of writing to that parity slot twice, we XOR the two data symbols together inside a SIMD register first and then write the merged result once. This halves the parity writes for those positions while keeping the read pattern localized like the horizontal traversal. As a result, Neighbor-Merging incurs fewer redundant memory accesses than both the horizontal and the vertical traversals in practice.
}

\subsection{Erasure-Parity Mapping}

\begin{frame}[t]{Key Idea 3 -- Erasure-Parity Mapping}
\footnotesize
\begin{columns}[T,totalwidth=\textwidth]
  \column{0.42\textwidth}
  \begin{itemize}
    \setlength{\itemsep}{0.45em}
    \item Map each erased chunk (data \emph{or} parity) to a survived parity for joint recovery.
    \item Decoded data symbols are \textbf{immediately} used to update erased parity --- no separate repair pass.
    \item One unified decoding flow for \emph{all} erasure patterns.
  \end{itemize}

  \vspace{0.4em}
  \begin{block}{Effect}
  \small Naive: 3 passes (substitute $\to$ eliminate $\to$ \textbf{repair}). \\
  FastTT: 2 passes; erased parity repaired \emph{along the way}.
  \end{block}

  \column{0.58\textwidth}
  \centering
  \includegraphics[width=\linewidth,height=0.64\textheight,keepaspectratio]{Erasure-Parity_Mapping_Decoding_Demonstration.pdf}

  \vspace{0.1em}
  {\scriptsize $(k,m)=(4,3)$. Lost: $\boldsymbol{d}_1,\boldsymbol{d}_2,\boldsymbol{c}_2$. Survived: $\boldsymbol{d}_3,\boldsymbol{d}_4,\boldsymbol{c}_1,\boldsymbol{c}_3$. \\ Dashed arrows $=$ operations \textbf{eliminated} by fusing repair into decoding.}
\end{columns}
\end{frame}
\note{
In the (4,3) example with $d_1$, $d_2$, $c_2$ lost: naive decoding substitutes $d_3$, $d_4$ into $c_1$, $c_3$, runs elimination to recover $d_1$ and $d_2$, then performs a separate third pass to re-encode $c_2$ from scratch. That third pass re-reads all data, which is redundant. FastTT maps each erased chunk to a survived parity: $d_1 \to c_1$, $d_2 \to c_3$, $c_2 \to c_2$ itself. During elimination, every decoded symbol of $d_1$ or $d_2$ is immediately XORed into the $c_2$ buffer. By the time elimination finishes, $c_2$ is already repaired. The dashed arrows in the figure show the operations eliminated by this fusion.
}

\subsection{Traversal Pattern Selection}

\begin{frame}[t]{Key Idea 4 -- Traversal Pattern Selection}
\footnotesize
\begin{columns}[T,totalwidth=\textwidth]
  \column{0.44\textwidth}
  \begin{itemize}
    \setlength{\itemsep}{0.35em}
    \item Only \textbf{1 chunk} lost $\Rightarrow$ each survived chunk read \textbf{once}, result written \textbf{once}.
    \item Vertical pattern: $\Theta((k{+}1)l)$. \\
          Horizontal: $\Theta((k{+}m)l)$ --- extra parity reads/writes for the $m{-}1$ untouched parities.
    \item No merging needed: each data chunk contributes exactly once, so pairwise XOR adds no benefit.
    \item All accesses are \textbf{sequential \& single-pass} $\Rightarrow$ use SIMD \textbf{streaming stores} (bypass cache, avoid pollution).
  \end{itemize}

  \vspace{0.3em}
  \begin{block}{Single-erasure vs. ISA-L}
  \scriptsize
  \textbf{+135\%\,--\,+182\%} improvement \\
  (i.e., FastTT is \textbf{2.35$\times$\,--\,2.82$\times$} the speed of ISA-L)
  \end{block}

  \column{0.56\textwidth}
  \centering
  \includegraphics[width=\linewidth,height=0.62\textheight,keepaspectratio]{Vertical_Traversal_Pattern_Demonstration.pdf}

  \vspace{0.1em}
  {\scriptsize $(k,m)=(4,3)$. Lost: $\boldsymbol{d}_2$ only. \\ Reconstructed by one XOR pass over $\boldsymbol{d}_1,\boldsymbol{d}_3,\boldsymbol{d}_4$ with shifted offsets from $\boldsymbol{c}_1$. \\ No elimination, no merging --- just stream-read, XOR, stream-write.}
\end{columns}
\end{frame}
\note{
When only one chunk is lost, each survived chunk is read exactly once and the result is written once. Vertical pattern costs $\Theta((k+1)l)$; horizontal would still touch the $m-1$ untouched parities, costing $\Theta((k+m)l)$. There is also no benefit from neighbor-merging: each data chunk contributes exactly one XOR, nothing to merge. Because all accesses are sequential and single-pass, we use SIMD streaming stores to bypass the cache entirely, avoiding pollution and maximizing bandwidth. In the (4,3) example with only $d_2$ lost, we stream-read $d_1$, $d_3$, $d_4$ with their shifted offsets from $c_1$, XOR them, and stream-write directly to $d_2$. No elimination needed. Result: single-erasure decoding is 135--182\% faster than ISA-L, i.e., 2.35--2.82$\times$ its speed. Multi-erasure falls back to EPM.
}

% =====================================================================
% Flagship results
% =====================================================================
\section{Flagship Results}
\subsection{Industrial and Academic Baselines}

\begin{frame}[t]{Flagship Results -- Throughput in Ceph}
\footnotesize
\begin{columns}[T,totalwidth=\textwidth]
  \column{0.5\textwidth}
  \begin{block}{vs. Ceph production baselines}
  \scriptsize
  Intel i9-14900K $\cdot$ AVX2 $\cdot$ 2\,MB chunks
  \begin{itemize}
    \item vs. ISA-L: \textbf{$+59.1\%$} encode, \textbf{$+50.9\%$} decode
    \item vs. Jerasure: \textbf{$+261\%$} encode, \textbf{$+161\%$} decode
  \end{itemize}

  \vspace{0.2em}
  \resizebox{\linewidth}{!}{%
  \begin{tabular}{lrrrrrr}
    \toprule
    & \multicolumn{3}{c}{Encode (GB/s)} & \multicolumn{3}{c}{Avg. decode} \\
    \cmidrule(lr){2-4} \cmidrule(lr){5-7}
    $(k,m)$ & \textbf{FastTT} & ISA-L & Jera. & \textbf{FastTT} & ISA-L & Jera. \\
    \midrule
    $(6,2)$ & 34.7 & 24.4 & 17.0 & 60.8 & 31.6 & 29.5 \\
    $(6,3)$ & 31.7 & 17.6 &  7.6 & 46.2 & 27.5 & 20.7 \\
    $(8,3)$ & 31.4 & 17.0 &  7.1 & 46.9 & 24.9 & 17.7 \\
    $(10,4)$ & 17.6 & 13.6 &  4.6 & 32.2 & 21.9 & 13.5 \\
    \bottomrule
  \end{tabular}}
  \end{block}

  \column{0.5\textwidth}
  \begin{block}{vs. academic SOTA (Cerasure)}
  \scriptsize
  \begin{itemize}
    \item \textbf{$+13.2\%$} encode on average
    \item \textbf{$+16.0\%$} decode on average; peak $+43.6\%$
    \item \textbf{Single-erasure: $+$135\%\,--\,$+$182\%} vs. ISA-L \\(2.35$\times$\,--\,2.82$\times$ the speed).
  \end{itemize}

  \vspace{0.2em}
  \resizebox{\linewidth}{!}{%
  \begin{tabular}{lrrrr}
    \toprule
    & \multicolumn{2}{c}{Encode (GB/s)} & \multicolumn{2}{c}{Avg. decode} \\
    \cmidrule(lr){2-3} \cmidrule(lr){4-5}
    $(k,m)$ & \textbf{FastTT} & Cerasure & \textbf{FastTT} & Cerasure \\
    \midrule
    $(6,2)$ & 48.0 & 44.2 & 67.6 & 59.7 \\
    $(6,3)$ & 32.0 & 31.0 & 49.4 & 44.8 \\
    $(8,3)$ & 32.8 & 28.7 & 48.7 & 43.9 \\
    $(10,4)$ & 24.6 & 19.4 & 37.3 & 33.4 \\
    \bottomrule
  \end{tabular}}
  \end{block}
\end{columns}

\vspace{0.2em}
{\tiny Avg. decode $=$ averaged over erasure counts $1..m$, for both data- and parity-chunk losses.}
\end{frame}
\note{
All three baselines --- ISA-L, Jerasure, and Cerasure --- are RS-based libraries. Shift-XOR codes have no prior production-grade implementations in storage systems, so this is also the first time Two-tone is benchmarked head-to-head against optimized RS implementations under the same Ceph framework.

Encoding-wise, FastTT consistently outperforms all baselines, including the academic state-of-the-art Cerasure, confirming that our memory-aware design effectively translates Two-tone's theoretical advantage into practical throughput.

For decoding, the average improvement is largely driven by the single-erasure case, where the TPS fast path uses streaming stores to deliver 2.35--2.82$\times$ the speed of ISA-L. As the number of erasures grows, the gap narrows: the single-erasure fast path no longer applies, and neighbor-merging cannot exploit the erased chunks during substitution because their content is unknown. Even so, FastTT still leads or ties the baselines on average across all configurations.
}

\subsection{Ablation of the Four Techniques}

\begin{frame}[t]{Flagship Results -- Ablation on $(8,3)$}
\footnotesize
\begin{columns}[T,totalwidth=\textwidth]
  \column{0.36\textwidth}
  \begin{itemize}
    \setlength{\itemsep}{0.45em}
    \item \textbf{EPM}: first jump in \emph{decode} \\ \small ($7.5 \to 16.9$, $+100\%$) -- fusion matters.
    \item \textbf{CFWP}: lifts both encode \& decode \\ \small ($+77\%$ enc, $+43\%$ dec).
    \item \textbf{NMA}: pure encoding gain \\ \small ($23.3 \to 31.2$, $+34\%$).
    \item \textbf{TPS}: second decode jump \\ \small ($24.1 \to 35.2$, $+46\%$).
  \end{itemize}

  \vspace{0.5em}
  \begin{block}{Takeaway}
  \small Each technique targets a distinct bottleneck; contributions are additive.
  \end{block}

  \column{0.64\textwidth}
  \centering
  \resizebox{\linewidth}{!}{%
  \begin{tabular}{ccccrr}
    \toprule
    \shortstack{Cache-Friendly\\Window Part.\\\textbf{(CFWP)}} &
    \shortstack{Neighbor-\\Merging Access\\\textbf{(NMA)}} &
    \shortstack{Erasure-Parity\\Mapping\\\textbf{(EPM)}} &
    \shortstack{Traversal\\Pattern Sel.\\\textbf{(TPS)}} &
    \textbf{Encode} & \textbf{Decode} \\
    \midrule
      &   &   &   & 13.2 &  7.5 \\
      &   & \checkmark &   & 13.2 & 16.9 \\
    \checkmark &   & \checkmark &   & 23.3 & 24.1 \\
    \checkmark & \checkmark & \checkmark &   & 31.2 & 24.1 \\
    \checkmark & \checkmark & \checkmark & \checkmark & \textbf{32.8} & \textbf{35.2} \\
    \bottomrule
  \end{tabular}}

  \vspace{0.4em}
  {\scriptsize Throughput in GB/s; decode averaged over erasure counts $1..m$.}
\end{columns}
\end{frame}
\note{
The ablation on (8,3) decomposes where the gains come from. EPM alone more than doubles decoding throughput, which validates our claim that fusing data recovery with parity repair removes a real bottleneck --- the redundant memory traffic of the separate repair stage. CFWP then improves both sides because cache-sized windows benefit any traversal, encoding or decoding. NMA is encoding-specific by design: it merges adjacent data contributions before writing parity, which only applies on the encoding side, so decoding stays flat here. Finally TPS gives a second large decoding jump by routing single-erasure cases to the streaming-store fast path. The four techniques target different bottlenecks, and the table shows their contributions are additive --- each one matters.
}

% =====================================================================
% Summary
% =====================================================================
\section{Summary and Future Research}
\subsection{Summary}

\begin{frame}[t]{Summary and Future Research}
\small
\begin{block}{One-sentence summary}
\design is the first high-performance systems implementation of \twotone \sxor coding, turning a theoretically efficient code family into a practical high-throughput storage primitive.
\end{block}

\vspace{0.4em}
\textbf{What to remember}
\begin{itemize}
  \setlength{\itemsep}{0.2em}
  \item Bottleneck $=$ \textbf{memory behavior}, not arithmetic.
  \item Four coordinated optimizations: \textbf{CFWP + NMA} (encode), \textbf{EPM + TPS} (decode).
  \item Single-erasure recovery benefits most -- the \textbf{common} case in production.
\end{itemize}

\vspace{0.4em}
\textbf{Open directions}
\begin{itemize}
  \setlength{\itemsep}{0.2em}
  \item Adaptive / heterogeneous window sizing beyond fixed $W$.
  \item Multi-level SIMD-aware merging beyond pairwise.
  \item Multi-erasure recovery: decompose into reusable single-erasure steps.
\end{itemize}
\end{frame}
\note{
FastTT is the first high-performance systems implementation of Two-tone Shift-XOR coding. The bottleneck is memory behavior, not arithmetic. Four coordinated techniques: CFWP and NMA for encoding, EPM and TPS for decoding. Single-erasure recovery, the dominant case in production, benefits most. Open directions: adaptive window sizing, multi-level SIMD merging, and decomposing multi-erasure patterns into reusable single-erasure steps. Thank you.
}

\subsection{Thanks}

\begin{frame}[plain]
  \centering
  \vfill
  {\Huge Thank you!}\\[1.2em]
  {\Large Questions \& Discussion}\\[2em]
  {\footnotesize Duo Sun, Ximing Fu, Qiang Wang \\[0.3em] HIT (Shenzhen) $\cdot$ Pengcheng Laboratory}
  \vfill
\end{frame}
\note{
Thank you for listening. I am happy to discuss the implementation details, the memory-access modeling behind each of the four techniques, the comparison with Reed-Solomon libraries, or the implications for production storage deployments in Ceph.
}

\end{document}
